






\documentclass[11pt,twoside,a4paper]{article}
%=============== En−Tete ===============
%−−− Insertion de paquetages (optionnel) −−−
\usepackage[french]{babel} % pour dire que le texte est en fran¸cais

\usepackage{a4} % pour la taille
\usepackage[T1]{fontenc} % pour les font postscript
\usepackage[cyr]{aeguill} % Police vectorielle TrueType, guillemets fran¸cais
\usepackage{epsfig} % pour g´erer les images
\usepackage{amsmath, amsthm} % tr`es bon mode math´ematique
\usepackage{amsfonts,amssymb}% permet la definition des ensembles
\usepackage{float} % pour le placement des figure
\usepackage{url} % pour une gestion efficace des url
\bibliographystyle{plain} % Style bibliographique
%=============== Corps ===============
\begin{document}


% Page de couverture
\begin{titlepage}
    \centering
    \vspace*{1cm}
    \includegraphics[width=4cm]{logo.jpg} % Assurez-vous d'avoir le fichier logo_universite.jpg dans le même répertoire que votre fichier .tex, ou indiquez le chemin complet.
    \vfill
    {\scshape\Large Université Des Antilles \par} % Assurez-vous de mettre en forme le nom de l'université selon vos préférences.
    \vspace{1cm}
    {\scshape\Large Master Mathématiques et Applications  \par} % Assurez-vous de mettre en forme le nom de la formation selon vos préférences.
    {\scshape\Large Economie et Finance  \par}
    \vspace{3cm}
    {\huge\bfseries Devoir d'économie\par}
    \vspace{1cm}
    {\Large\itshape Pierre Chrislin DORIVAL\par} % Assurez-vous de mettre en forme votre nom et prénom selon vos préférences.
    \vspace{0.5cm}
    {\Large\itshape Professeur: Patrice BORDA\par} % Assurez-vous de mettre en forme le nom du professeur selon vos préférences.
    \vfill
    {\large 25 Avril  2024 \par} % Assurez-vous de mettre la date selon vos préférences.
\end{titlepage}

% Page de formules mathématiques avec démonstration du théorème de Cauchy
\section*{Résoluton Exercices 1}

On considère un modèle à générations imbriquées où les ménages ne vivent
que deux périodes. Au cours de leur période d'activité, ils perçoivent un
salaire $w_t$ qu'ils affectent à la consommation $c_t$ courante et à l'épargne $s_t$.
Au cours de leur période de retraite, ils consacrent leur épargne à la consom-
mation différée $d_{t+1}$. Leurs préférences intertemporelles sont représentées par
la fonction d'utilité :


$u_t(c_t; d_{t+1})=\frac{1}{2}\log{c_t} + \frac{1}{2}\log{d_{t+1}}$\\ \\


Dans cette économie, les impôts sont nuls $(x = 0)$. Les individus nés en $t$
sont en nombre $L_t$. Le taux de croissance du nombre de travailleurs d'une
génération à l'autre est noté $n = 1$:\\ \\
$a)$ Exprimons les consommations des jeunes et des vieux en fonction du saliare
$wt$:\\

Fonction d'utilité compte tenu de la contrainte budgétaire intertemporelle:

$ \left\{\begin{array}{rl}
 Max \ u_t(c_t; d_{t+1}) &=  \frac{1}{2}\log{c_t} + \frac{1}{2}\log{d_{t+1}} \\
 
w_t & =  c_t + \frac{d_{t+1}}{1+r_{t+1}}\end{array}\right. $  \\ \\
Considérons la fonction Lagrangienne :\\

$ L(c_t,d_{t+1},\lambda_t)= \frac{1}{2}\log{c_t} + \frac{1}{2}\log{d_{t+1}} +\lambda_t(w_t-c_t - \frac{d_{t+1}}{1+r_{t+1}})  $  \\

Faisons la dérivée de la fonction Lagrangienne par rapport $c_t :  $ \\

$
\frac{\partial L(c_t,d_{t+1},\lambda_t)}{\partial c_t} = \frac{1}{2c_t} - \lambda_t 
$ \\ \\

Posons  \ $ \frac{\partial L(c_t,d_{t+1},\lambda_t)}{\partial c_t} = 0  $ \\



\begin{equation}\label{eq:test1}
  \Rightarrow \lambda_t=\frac{1}{2c_t}
\end{equation} 


Faisons la dérivée de la fonction Lagrangienne par rapport $d_{t+1} : $\\ 

$
\frac{\partial L(c_t,d_{t+1},\lambda_t)}{\partial d_{t+1}} = \frac{1}{2 d_{t+1}} - \frac{\lambda_t}{1+r_{t+1}} 
$\\ \\

Posons \ $ \frac{\partial L(c_t,d_{t+1},\lambda_t)}{\partial d_{t+1}} = 0 $ \\

  \begin{equation}\label{eq:test2}
  \Rightarrow \lambda_t=\frac{1+r_{t+1}}{2 d_{t+1}}
\end{equation}  

Faisons la dérivée de la fonction Lagrangienne par rapport $\lambda_t :  $ \\ \\
 
$
\frac{\partial L(c_t,d_{t+1},\lambda_t)}{\partial \lambda_t} = w_t-c_t - \frac{d_{t+1}}{1+r_{t+1}} 
$ \\

Posons \ $ \frac{\partial L(c_t,d_{t+1},\lambda_t)}{\partial \lambda_t} = 0 $ \\

 \begin{equation}\label{eq:test3}
  \Rightarrow w_t=c_t + \frac{d_{t+1}}{1+r_{t+1}}
\end{equation}
    

(\ref{eq:test1}) et (\ref{eq:test2}) nous donnent : \\ \\

$
\frac{1}{c_t} = \frac{1+r_{t+1}}{d_{t+1}}
$\\ \\

$ \Rightarrow c_t= \frac{d_{t+1}}{1+r_{t+1}}$  et  $ d_{t+1}= c_t ({1+r_{t+1}})$\\

Dans (\ref{eq:test1})\: $ w_t=c_t + \frac{d_{t+1}}{1+r_{t+1}}$ \\

D'o\'{u} \:
 \fbox {$c_t = \frac{1}{2} w_t $} \: et \: \fbox {$d_{t+1} = \frac{1}{2} w_t(1+r_{t+1}) $}\\ \\
 $b)$  Déterminer l'épargne en fonction de $w_t$.\\
 
 On sait que: $ s_t = \frac{d_{t+1}}{1+r_{t+1}} $\\
 
 $ \Rightarrow$  \fbox {$s_t = \frac{1}{2} w_t $} \\
 
 La technologie est représentée par la fonction de production :\\
 
 $Y_t=10K_t^{\frac{1}{2}}L_t^{\frac{1}{2}}$\\
 
 avec $K_t$ est le stock de capital et $L_t$ l'emploi.\\ \\
 $c)$ Supposant les marchés concurrentiels, exprimons le salaire et les taux d'intér\^{e}t en fonction du stock de capital par t\^{e}te.\\

$ max\ \pi_t= Y_t -w_tL_t-K_t r_t $\\

 $ max\ \pi_t=10K_t^{\frac{1}{2}}L_t^{\frac{1}{2}}-w_tL_t-K_t r_t $\\
 
 $
  \frac{\partial (max\ \pi_t)}{\partial K_t}=5K_t^{\frac{-1}{2}}L_t^{\frac{1}{2}}-r_t $\\
 
$ \Rightarrow r_t= 5(\frac{K_t}{L_t})^{\frac{-1}{2}}=5k_t^{\frac{-1}{2}}$ \: avec $k_t$ stock capital par t\^{e}te.\\


 $ \Rightarrow$  \fbox {$r_t = 5k_t^{\frac{-1}{2}} $} \\
 \\Calcul maintenant $\frac{\partial (max\ \pi_t)}{\partial L_t}$\\
 
 $\frac{\partial (max\ \pi_t)}{\partial L_t}=5K_t^{\frac{1}{2}}L_t^{\frac{-1}{2}}-w_t$\\
 
 $ \Rightarrow r_t= 5(\frac{K_t}{L_t})^{\frac{1}{2}}=5k_t^{\frac{1}{2}}$ \: avec $k_t$ stock capital par t\^{e}te.\\

$ \Rightarrow$  \fbox {$w_t = 5k_t^{\frac{1}{2}} $} \\ \\
$d)$ Déterminons l'évolution du stock de capital par tête.\\\\
$K_{t+1}=L_t s_t$ \: or \: $L_t=\frac{L_{t+1}}{N}$\\\\
$\Rightarrow  K_{t+1}=L_{t+1}\frac{s_t}{2}$\\\\
$\Rightarrow \frac{K_{t+1}}{L_{t+1}}=k_{t+1}=\frac{1}{4}w_t$ \: car $s_t = \frac{1}{2} w_t $\\\\
$ \Rightarrow$  \fbox {$k_{t+1} =\frac{5}{4}k_t^{\frac{1}{2}} $} \: car \ $w_t = 5k_t^{\frac{1}{2}} $\\\\
$e)$ Déterminons l'expression du capital par tête à l'équilibre. Cet équilibre
est-il stable:\\\\

A l'équilibre: \: $k_{t+1}=k_t=k^*=\frac{5}{4}k^{*\frac{1}{2}}$\\

$ \Rightarrow$  \fbox {$k^*=\frac{25}{16}$} \\

$\log{(k_{t+1})}=\log{(\frac{5}{4}k_t^{\frac{1}{2}})}+\log{(\frac{5}{4})}+\frac{1}{2}\log{(k_t)}$\\

$
\Rightarrow \log{(k_{t+1})} -\log{(\frac{5}{4})}=\frac{1}{2}\log{(k_t)}$\\

$
\Rightarrow \log{(\widehat{k_{t+1}})}=\frac{1}{2}\log{(\widehat{k_{t}})}$\\

$
 \alpha <\frac{1}{2}<1$\\
 L'équilibre est stable\\\\
 $ f)$ La valeur du stock de capital de la règle d'or.\\\\
 $Y_t=10K_t^{\frac{1}{2}}L_t^{\frac{1}{2}}$\\\\
 $Y_t=L_tc_t+L_{t+1}K_{t+1}+L_{t-1}d_{t+1}$\\\\
 $\frac{Y_t}{L_t}=c_t+(1+n)K_{t+1}+\frac{d_{t+1}}{1+n}$\: car $L_{t+1}=(1+N)L_t$\\\\     
$ \frac{Y_t}{L_t}=10(\frac{K_t}{L_t})^{\frac{1}{2}}=10k_t^{\frac{1}{2}}$\\

$ 10k_t^{\frac{1}{2}}=c_t+(1+n)K_{t+1}+\frac{d_{t+1}}{1+n}$\\
 
 $ c_t=10k_t^{\frac{1}{2}}-2K_{t+1}-\frac{d_{t+1}}{2}$\\
 
 $\frac{\partial c_t}{K_t}=5k_t^\frac{-1}{2}-2$\\
 
  Posons $\frac{\partial c_t}{K_t}=0$\\
  
  $k_t^\frac{-1}{2}=\frac{2}{5}$\\
  
  $k_t=\frac{25}{4}=k_{or}$\\

  $k^*<k_{or}$\\\\
  Dans ce cas l'économie est en régime sous-accumulation.\\
  
  \section*{Résoluton Exercices 2}
  
  Considérons un modèle à générations imbriquées où les ménages ne vivent que deux périodes. Durant la première période, ils perçoivent un salaire $w_t$; qu'ils affectent à la  consommation courante $c_t$; et à l'épargne $s_t$; et au versement des cotisations de retraite, $\sigma w_t$: $\sigma$; est le taux de cotisation à la retraite par répartition. Devenus vieux à la deuxième période, ils consacrent leur épargne et leur pension de retraite $\overline{\theta}\phi w_t$; à la consommation différée $d_{t+1}$.
$\overline{\theta}$ ; est l'espérance de vie du retraité tandis que ; est le taux de remplacement
des retraites. Les préférences des ménages sont représentées par la fonction d'utilité :\\
$u_t(c_t; d_{t+1})=(1-s)\log{c_t} + s\log{d_{t+1}}$ \: $0<s<1$
\\
Notons $N$, l'effctif (constant) de la population active.\\

$a)$ Déterminons la contrainte budgétaire intertemporelle des ménages.\\

Fonction d'utilité compte tenu de la contrainte budgétaire intertemporelle:\\

$ \left\{\begin{array}{rl}
 Max \ u_t(c_t; d_{t+1})=  (1-s)\log{c_t} + s\log{d_{t+1}} ; 0<s<1 \\
 
w_t  =  c_t + s_t +  \sigma w_t \end{array}\right. $  \\

On sait que :\\

$d_{t+1}=(1+r_{t+1})s_t+\overline{\theta}\phi w_t$\\

  $
\Rightarrow s_t=\frac{d_{t+1}-\overline{\theta}\phi w_t}{(1+r_{t+1})}$\\

%D'o\'{u}:\\
%$\Rightarrow w_t  -\sigma w_t =  c_t + \frac{d_{t+1}-\overline{\theta}\phi w_t}{(1+r_{t+1})}$\\ 
%$\Rightarrow w_t = \frac{c_t + \frac{d_{t+1}-\overline{\theta}\phi w_t}{(1+r_{t+1})}}{(1-\sigma)} $\\
%D'o\'{u}, la fonction d'utilité compte tenu de la contrainte budgétaire intertemporelle:\\
%$ \left\{\begin{array}{rl}
% Max \ u_t(c_t; d_{t+1})=  (1-s)\log{c_t} + s\log{d_{t+1}} ; 0<s<1 \\
 
%\Rightarrow w_t = \frac{c_t(1+r_{t+1}) + d_{t+1}-\overline{\theta}\phi w_t}{(1-\sigma)((1+r_{t+1}))} \end{array}\right. $  \\
 
 b) Exprimeons les consommations et l'épargne en fonction du salaire et du taux
d'intér\^{e}t.\\
La technologie des firmes individuelles est représentée par la fonction de production suivante :\\
$Y_{it}=A_tK^{\alpha}_{it}L^{1-\alpha}_{it}$\\

avec \: $A_t=BK^{1-\alpha}_t$ \: $0<\alpha<1$ \: $ i=\{ 1,2,3,...,M \}$\\

$Y_{it}, K_{it},L_{it},$ ; reprÈsentent respectivement, la production, le stock de capital et l'emploi de la firme $i$. $A_t$ est l'efficacité du travail tandis que $K_t$ est assimilé au stock de capital total.\\



Considérons la fonction Lagrangienne :

$ L(c_t,d_{t+1},\lambda_t)= (1-s)\log{c_t} + s\log{d_{t+1}} +\lambda_t(w_t-c_t - \sigma w_t-\frac{d_{t+1}-\overline{\theta}\phi w_t}{(1+r_{t+1})})  $  \\

Faisons la dérivée de la fonction Lagrangienne par rapport $c_t :  $ \\

$\\
\frac{\partial L(c_t,d_{t+1},\lambda_t)}{\partial c_t} = \frac{1-s}{c_t} - \lambda_t 
$ \\
Posons  $ \frac{\partial L(c_t,d_{t+1},\lambda_t)}{\partial c_t} = 0  $ \\



\begin{equation}\label{eq:test4}
  \Rightarrow \lambda_t=\frac{1-s}{c_t}
\end{equation} 


Faisons la dérivée de la fonction Lagrangienne par rapport $d_{t+1} : \\ $
$\\
\frac{\partial L(c_t,d_{t+1},\lambda_t)}{\partial d_{t+1}} = \frac{s}{ d_{t+1}} - \frac{\lambda_t}{1+r_{t+1}} 
$\\

Posons  $ \frac{\partial L(c_t,d_{t+1},\lambda_t)}{\partial d_{t+1}} = 0 $ \\



  \begin{equation}\label{eq:test5}
  \Rightarrow \lambda_t=\frac{(1+r_{t+1})s}{ d_{t+1}}
\end{equation}  

Faisons la dérivée de la fonction Lagrangienne par rapport $\lambda_t :  $ \\
$\\
\frac{\partial L(c_t,d_{t+1},\lambda_t)}{\partial \lambda_t} = w_t-c_t - \frac{d_{t+1}}{1+r_{t+1}} 
$ \\ \\
Posons  $ \frac{\partial L(c_t,d_{t+1},\lambda_t)}{\partial \lambda_t} = 0 $ \\ \\
$\Rightarrow w_t-c_t - \sigma w_t-\frac{d_{t+1}-\overline{\theta}\phi w_t}{(1+r_{t+1})}=0$\\ \\
 \begin{equation}\label{eq:test6}
  \Rightarrow w_t = \frac{c_t(1+r_{t+1}) + d_{t+1}-\overline{\theta}\phi w_t}{(1-\sigma)(1+r_{t+1})}
\end{equation}
    

(\ref{eq:test4}) et (\ref{eq:test5}) nous donnent : \\

$\\
\frac{1-s}{c_t} = \frac{(1+r_{t+1})s}{ d_{t+1}}\\
$\\
$ \Rightarrow c_t= \frac{(1-s)d_{t+1}}{(1+r_{t+1})s}$  \: et \: $ d_{t+1}=c_t \frac{ s{(1+r_{t+1})}}{1-s}$\\ \\
Dans (\ref{eq:test6})\: $ w_t = \frac{c_t(1+r_{t+1}) + d_{t+1}-\overline{\theta}\phi w_t}{(1-\sigma)((1+r_{t+1}))}$\\ \\
Rempla\c{c}ons $d_{t+1}$ dans $w_t$, \\ \\
$c_t(1+r_{t+1})+\frac{c_t s{(1+r_{t+1})}}{1-s}-\overline{\theta}\phi w_t= w_t(1_\sigma)(1+r_{t+1})$ \\ \\
D'o\'{u} \:
 \fbox {$c_t =  w_t(1-s)(\frac{\overline{\theta}\phi}{1+r_{t+1}} +1- \sigma )$} \\ \\ 
 Portons $c_t$ dans $d_{t+1}$, \\ \\
 $ d_{t+1}= w_t(1-s)(\frac{\overline{\theta}\phi}{1+r_{t+1}} +1- \sigma )\frac{ s{(1+r_{t+1})}}{1-s}$ \\ \\
\fbox {$d_{t+1} = s w_t\left(\overline{\theta}\phi +(1- \sigma)(1+r_{t+1}) \right)$}\\ \\
 On sait que :  \\ \\
 $s_t=\frac{d_{t+1}-\overline{\theta}\phi w_t}{(1+r_{t+1})}$\\ \\ 
 Rempla\c{c}ons $d_{t+1}$ dans $w_t$, \\ \\
  $s_t=\frac{s w_t\left(\overline{\theta}\phi +(1- \sigma)(1+r_{t+1}) \right)-\overline{\theta}\phi w_t}{(1+r_{t+1})}$\\ \\
\fbox {$s_t = w_t\left(\frac{\overline{\theta}\phi}{1+r_{t+1}}(s-1) +(1- \sigma)s \right)$}\\ \\
$c)$ Montrons que dans une économie concurrentielle, le salaire et le taux d'intér\^{e}t sont donnés par :\\ \\
$w_t=(1-\alpha)BN^{1-\alpha}k_t$\\ \\
$r_t=\alpha BN^{1-\alpha}-1$\\ \\
avec \ avec $k_t=\frac{K_t}{N}$\\ \\
Considérons  \ $ max\ \pi_{it}= Y_{it} -w_tL_{it}-K_{it} r_ti-\sigma K_{it} $ \ avec $\sigma=1$\\ \\
\ $ max\ \pi_{it}= A_tK^{\alpha}_{it}L^{1-\alpha}_{it} -w_tL_{it}-K_{it} r_{ti} -K_{it} $ \\ \\
 $ \frac{\partial (max\ \pi_t)}{\partial K_{it}}=\alpha A_tK^{\alpha-1}_{it}L^{1-\alpha}_{it}-r_t-1 = 0 $\\ \\
 $ \Rightarrow r_t=\alpha A_tK^{\alpha-1}_{it}L^{1-\alpha}_{it} -1$\\ \\
 $ \Rightarrow r_t=\alpha A_t \left(\frac{K_{it}}{L_{it}}\right)^{\alpha-1} -1$\\ \\ or \ $k_t=\frac{K_t}{N}$\\ \\
 $ \Rightarrow r_t=\alpha A_t k_t^{\alpha-1} -1$\\ \\ 
 $ \Rightarrow r_t=\alpha A_t \left(\frac{K_t}{N}\right)^{\alpha-1} -1$\\ \\ 
 $ \Rightarrow r_t=\alpha A_t K_t^{\alpha-1}N^{1-\alpha} -1$\\ \\ 
$ \Rightarrow r_t=\alpha BK^{1-\alpha}_t K_t^{\alpha-1}N^{1-\alpha} -1$\\ \\ 
Car : \ $A_t=BK^{1-\alpha}_t$\\ \\
 Par conséquent : \ \fbox{$r_t=\alpha BN^{1-\alpha} -1$}\\ \\ 
$\frac{\partial\ max\ \pi_{it} }{\partial L_{it}}=(1-\alpha)A_tK_{it}^{\alpha}L_{it}^{-\alpha}-w_t=0$\\ \\
$ \Rightarrow w_t=(1-\alpha)A_t\left(\frac{K_{it}}{L_{it}}\right)^{\alpha}$\\ \\
$ \Rightarrow w_t=(1-\alpha)A_tk_t^{\alpha}$\\ \\
$ \Rightarrow w_t=(1-\alpha)A_t\left(\frac{K_t}{N}\right)^{\alpha}$\\ \\
$ \Rightarrow w_t=(1-\alpha)BK^{1-\alpha}_t K_t^{\alpha}N^{-\alpha}$\\ \\
Car : \ $A_t=BK^{1-\alpha}_t$\\ \\
$ \Rightarrow w_t=(1-\alpha)B K_tN^{-\alpha}$\\ \\
Or \ $k_t=\frac{K_t}{N} \ \Rightarrow \ K_t=k_tN$\\ \\
$ \Rightarrow w_t=(1-\alpha)B k_tNN^{-\alpha}$\\ \\
Par conséquent : \ \fbox{$ w_t=(1-\alpha)B N^{1-\alpha}k_t$}\\ \\

$d)$ Déterminer la dynamique du stock de capital par t\^{e}te.\\ \\
$K_{t+1}=L_ts_t=Ns_t$\\ \\
$\Rightarrow K_{t+1}=Nw_t\left(\frac{\overline{\theta}\phi}{1+r_{t+1}}(s-1) +(1- \sigma)s \right)$\\ \\
$\Rightarrow K_{t+1}=N(1-\alpha)B N^{1-\alpha}k_t\left(\frac{\overline{\theta}\phi}{1+r_{t+1}}(s-1) +(1- \sigma)s \right)$\\ \\
$\Rightarrow \frac{K_{t+1}}{N}=(1-\alpha)B N^{1-\alpha}k_t\left(\frac{\overline{\theta}\phi}{1+r_{t+1}}(s-1) +(1- \sigma)s \right)$\\ \\
Par conséquent : \ \fbox{$k_{t+1}=(1-\alpha)B N^{1-\alpha}\left(\frac{\overline{\theta}\phi}{1+r_{t+1}}(s-1) +(1- \sigma)s \right )k_t$}\\ \\
$e)$ Sous quelle condition le stock de capital par tÍte d'équilibre est-il localement stable ?


$f)$ En posant  $(1+g_t)=\frac{k_{t+1}}{k_t}$ déterminer le taux de croissance de l'économie.\\ \\
$k_{t+1}=(1+g_t)k_t$\\
Or \ $k_{t+1}=(1-\alpha)B N^{1-\alpha}\left(\frac{\overline{\theta}\phi}{1+r_{t+1}}(s-1) +(1- \sigma)sk_t \right)$\\ \\
$(1-\alpha)B N^{1-\alpha}\left(\frac{\overline{\theta}\phi}{1+r_{t+1}}(s-1) +(1- \sigma)s \right)k_t=(1+g_t)k_t$\\ \\
$\Rightarrow g_t=(1-\alpha)B N^{1-\alpha}\left(\frac{\overline{\theta}\phi}{1+r_{t+1}}(s-1) +(1- \sigma)s \right)-1$\\ \\
Posons : $g_t=g^*>0$\\ \\
$\Rightarrow (1-\alpha)B N^{1-\alpha}\left(\frac{\overline{\theta}\phi}{1+r_{t+1}}(s-1) +(1- \sigma)s \right)>1$\\ \\


$g)$ Etudier analytiquement l'effet de l'allongement de l'espérance de vie des
retraités sur la dynamique de l'économie. Commenter

\end{document}






